% Source: https://artofproblemsolving.com/community/c3h2485993_some_problems
% Stable identity: aops-topic-2485993-q04
Let $\sigma(m)$ be the sum of all positive divisors of $m$; then $m$ is abundant exactly when $\sigma(m)>2m$. Check $n=1,2,\dots$:
\begin{itemize}
\item $n=1$: $\sigma(1)=1<2$;
\item $n=2$: $\sigma(4)=1+2+4=7<8$;
\item $n=3$: $\sigma(9)=1+3+9=13<18$;
\item $n=4$: $\sigma(16)=31<32$;
\item $n=5$: $\sigma(25)=31<50$;
\item $n=6$: $\sigma(36)=1+2+3+4+6+9+12+18+36=91>72$.
\end{itemize}
So the least such $n$ is $\boxed{6}$; indeed the proper divisors of $36$ sum to $55>36$.

\emph{Remark.} No prime power $n=p^k$ with $k\ge1$ works: $\dfrac{\sigma\left(p^{2k}\right)}{p^{2k}}=1+\dfrac1p+\cdots+\dfrac{1}{p^{2k}}<\dfrac{p}{p-1}\le2$. That is why $2,3,4,5$ fail, and $6$ is the smallest integer greater than $1$ that is not a prime power.
